\documentclass[11pt]{article}
\usepackage[margin=1in]{geometry}
\usepackage[T1]{fontenc}
\usepackage{lmodern}
\usepackage{amsmath,amssymb}
\usepackage[hidelinks]{hyperref}
\usepackage{microtype}
\raggedright

\begin{document}

\begin{center}
{\Large Literature answer for arXiv:0909.0131}\\[4pt]
{\large Every bounded truncated Toeplitz operator has a
Carleson-measure quasisymbol}
\end{center}

\paragraph{Status.}
\textbf{literature\_already\_answered}.  This is an exact later theorem,
not a new proof produced by the run.

\paragraph{Source question.}
Remark~5.3 of Baranov--Chalendar--Fricain--Mashreghi--Timotin
\cite{BCFMT}, source PDF page 19, asks whether every bounded truncated
Toeplitz operator in $\mathcal T(K_\Theta)$ is of the form
$A_\mu^\Theta$ for some complex Carleson measure $\mu$ for $K_\Theta$.
Here this means that $|\mu|$ induces a bounded embedding
$K_\Theta\hookrightarrow L^2(|\mu|)$ and
\[
 \langle A_\mu^\Theta f,g\rangle=\int_{\mathbb T}f\overline g\,d\mu.
\]

\paragraph{Supporting answer.}
Baranov--Bessonov--Kapustin \cite{BBK} define $\mathcal C_2(\theta)$ as
the finite complex Borel measures $\mu$ for which
$K_\theta\hookrightarrow L^2(|\mu|)$ is bounded, and call such a measure a
quasisymbol when the displayed identity holds.  Their introduction
explicitly identifies as Sarason's conjecture the assertion that every
bounded truncated Toeplitz operator has such a quasisymbol and says that they
prove it.  Theorem~2.1(2), supporting PDF page 6, states:
\begin{quote}
Any bounded truncated Toeplitz operator on $K_\theta$ admits a quasisymbol
from $\mathcal C_2(\theta)$.
\end{quote}
This is exactly the affirmative answer requested in the source.

\paragraph{Scope.}
The result answers the Carleson-measure representation question only.  It
does not settle the source's general reproducing-kernel thesis (Question~3)
or the later conjecture characterizing the inner functions for which every
bounded truncated Toeplitz operator has a bounded $L^\infty$ symbol.

\paragraph{Search evidence.}
Cheap run indexes and bounded web/arXiv searches through 2026-08-09 used
arXiv:0909.0131, the exact question, ``quasisymbol'', ``Carleson measure'',
and truncated Toeplitz symbols.  They located the explicit later theorem
arXiv:1009.5123 and corroborating surveys.

\begin{thebibliography}{9}
\bibitem{BCFMT}
A. Baranov, I. Chalendar, E. Fricain, J. Mashreghi, and D. Timotin,
\emph{Bounded symbols and reproducing kernel thesis for truncated Toeplitz
operators}, J. Funct. Anal. 259 (2010), 2673--2701;
arXiv:0909.0131.

\bibitem{BBK}
A. Baranov, R. Bessonov, and V. Kapustin,
\emph{Symbols of truncated Toeplitz operators},
J. Funct. Anal. 261 (2011), 3437--3456;
arXiv:1009.5123, doi:10.1016/j.jfa.2011.08.005.
\end{thebibliography}

\end{document}
